(1) 无源均匀介质中的电场满足 $\nabla^2\boldsymbol E+\omega^2\mu\epsilon\boldsymbol E=0$。理想金属壁上切向电场必须为零，故 $\sin(\pi x/a)$ 在两壁满足条件。将给定模态代入电场 Helmholtz 方程：
\[
-\left[(\pi/a)^2+k_z^2\right]E_y+\omega^2\mu\epsilon E_y=0,
\]
因此
\[
k_z^2=\mu\epsilon\omega^2-(\pi/a)^2
=\mu\epsilon(\omega^2-\omega_c^2),\quad
\omega_c=\frac{\pi}{a\sqrt{\mu\epsilon}}.
\]
$\omega>\omega_c$ 时选择 $k_z>0$ 描述向右传播模式，$\omega<\omega_c$ 时 $k_z$ 为纯虚数而非传播波数。

(2) 记介质均匀波速 $c_m=1/\sqrt{\mu\epsilon}$，高于截止时有
\[
k_z=\frac1{c_m}\sqrt{\omega^2-\omega_c^2},
\quad
v_{\rm ph}=\frac{\omega}{k_z}
=\frac{c_m}{\sqrt{1-\omega_c^2/\omega^2}},
\]
\[
v_g=\frac{d\omega}{dk_z}
=c_m\sqrt{1-\omega_c^2/\omega^2},\qquad
v_{\rm ph}v_g=c_m^2.
\]
从磁场横向分量可定义 TE 波阻抗 $Z_{\rm TE}=|E_y/H_x|=\omega\mu/k_z$。靠近截止时 $v_{\rm ph}\to\infty$，但能量与窄带波包传播速度 $v_g\to0$，所以相速度超过 $c_m$ 并不表示信息超光速。

(3) 在给定时间约定下，Faraday 定律为 $\nabla\times\boldsymbol E=i\omega\mu\boldsymbol H$。由于 $\boldsymbol E=E_y\hat y$，有 $(\nabla\times\boldsymbol E)_x=-\partial_zE_y$、$(\nabla\times\boldsymbol E)_z=\partial_xE_y$，因此
\[
H_x=-\frac{k_z}{\omega\mu}E_y,\qquad
H_z=-i\frac{\pi}{\omega\mu a}E_0\cos(\pi x/a)e^{i(k_z z-\omega t)}.
\]
平均坡印廷流的纵向分量是 $\overline S_z=-\frac12\Re(E_yH_x^*)$，故在 $k_z>0$ 的传播模式中
\[
\overline S_z(x)=\frac{k_z|E_0|^2}{2\omega\mu}\sin^2(\pi x/a).
\]
对每单位 $y$ 向宽度积分得到
\[
\overline P_{\rm per\,width}
=\int_0^a\overline S_z\,dx
=\frac{ak_z|E_0|^2}{4\omega\mu}>0.
\]

(4) 当 $0<\omega<\omega_c$，定义 $\kappa=\sqrt{(\pi/a)^2-\mu\epsilon\omega^2}>0$，选择沿 $+z$ 不发散的支为 $k_z=i\kappa$，电场幅值 $\propto e^{-\kappa z}$，所以 $1/e$ 振幅衰减长度为 $\ell_{\rm ev}=1/\kappa$。此时 $H_x=-i\kappa E_y/(\omega\mu)$，它相对于 $E_y$ 相差 $\pi/2$，因而 $\Re(E_yH_x^*)=0$，沿波导的周期平均功率流为零。这表示纯倏逝模式不向无限远传输功率，而非局部电磁场严格为零。
